\documentclass[11pt]{article}
\newcommand{\videoid}{qFL_OneuFwI}
% Preámbulo compartido de las notas del curso de topología algebraica.
% Cada nota empieza con \documentclass[11pt]{article}, define \videoid{...}
% (el identificador del video de YouTube) y luego hace % Preámbulo compartido de las notas del curso de topología algebraica.
% Cada nota empieza con \documentclass[11pt]{article}, define \videoid{...}
% (el identificador del video de YouTube) y luego hace % Preámbulo compartido de las notas del curso de topología algebraica.
% Cada nota empieza con \documentclass[11pt]{article}, define \videoid{...}
% (el identificador del video de YouTube) y luego hace \input{preambulo}.
% Compilar con LuaLaTeX:  lualatex nombre-de-la-nota.tex
\usepackage[spanish,provide=*]{babel}
\usepackage[T1]{fontenc}
\usepackage{lmodern}
\usepackage[a4paper,margin=2.6cm]{geometry}
\usepackage{amsmath,amssymb,amsthm}
\usepackage{xcolor}
\usepackage{tikz}
\usetikzlibrary{arrows.meta,babel}
\usepackage[colorlinks=true,urlcolor=blue!50!black,linkcolor=blue!50!black]{hyperref}

\hyphenpenalty=10000 \exhyphenpenalty=10000 \emergencystretch=3em

% Colores de las figuras (los mismos que usa la página)
\definecolor{extremo}{RGB}{47,107,79}
\definecolor{azul}{RGB}{44,74,154}

% Definiciones numeradas por unidad
\theoremstyle{definition}
\newtheorem{defn}{Definición}
% Numeración por unidad: \unidad{u}{n} (después de \input{preambulo}) hace que las
% definiciones de la nota se llamen u.n, u.(n+1), ...; n es el primer número
% libre de esa unidad (ver la tabla de CONVENCIONES.md).
\newcommand{\unidad}[2]{\renewcommand{\thedefn}{#1.\arabic{defn}}\setcounter{defn}{\numexpr#2-1\relax}}

% Observaciones y erratas del video. Usar SOLO si hay algo que señalar:
%   \begin{erratas}
%     \errata{203}{3:23}{Se dice ..., pero debe decir ...}
%   \end{erratas}
% El primer argumento son los segundos (enlaza al minuto exacto del video),
% el segundo es la marca visible, el tercero es el texto.
\newenvironment{erratas}{\section*{Observaciones y erratas del video}\begin{itemize}}{\end{itemize}}
\newcommand{\errata}[3]{%
  \edef\errurl{https://www.youtube.com/watch?v=\videoid\string&t=#1}%
  \item[\textbf{\expandafter\href\expandafter{\errurl}{#2}}] #3}
.
% Compilar con LuaLaTeX:  lualatex nombre-de-la-nota.tex
\usepackage[spanish,provide=*]{babel}
\usepackage[T1]{fontenc}
\usepackage{lmodern}
\usepackage[a4paper,margin=2.6cm]{geometry}
\usepackage{amsmath,amssymb,amsthm}
\usepackage{xcolor}
\usepackage{tikz}
\usetikzlibrary{arrows.meta,babel}
\usepackage[colorlinks=true,urlcolor=blue!50!black,linkcolor=blue!50!black]{hyperref}

\hyphenpenalty=10000 \exhyphenpenalty=10000 \emergencystretch=3em

% Colores de las figuras (los mismos que usa la página)
\definecolor{extremo}{RGB}{47,107,79}
\definecolor{azul}{RGB}{44,74,154}

% Definiciones numeradas por unidad
\theoremstyle{definition}
\newtheorem{defn}{Definición}
% Numeración por unidad: \unidad{u}{n} (después de % Preámbulo compartido de las notas del curso de topología algebraica.
% Cada nota empieza con \documentclass[11pt]{article}, define \videoid{...}
% (el identificador del video de YouTube) y luego hace \input{preambulo}.
% Compilar con LuaLaTeX:  lualatex nombre-de-la-nota.tex
\usepackage[spanish,provide=*]{babel}
\usepackage[T1]{fontenc}
\usepackage{lmodern}
\usepackage[a4paper,margin=2.6cm]{geometry}
\usepackage{amsmath,amssymb,amsthm}
\usepackage{xcolor}
\usepackage{tikz}
\usetikzlibrary{arrows.meta,babel}
\usepackage[colorlinks=true,urlcolor=blue!50!black,linkcolor=blue!50!black]{hyperref}

\hyphenpenalty=10000 \exhyphenpenalty=10000 \emergencystretch=3em

% Colores de las figuras (los mismos que usa la página)
\definecolor{extremo}{RGB}{47,107,79}
\definecolor{azul}{RGB}{44,74,154}

% Definiciones numeradas por unidad
\theoremstyle{definition}
\newtheorem{defn}{Definición}
% Numeración por unidad: \unidad{u}{n} (después de \input{preambulo}) hace que las
% definiciones de la nota se llamen u.n, u.(n+1), ...; n es el primer número
% libre de esa unidad (ver la tabla de CONVENCIONES.md).
\newcommand{\unidad}[2]{\renewcommand{\thedefn}{#1.\arabic{defn}}\setcounter{defn}{\numexpr#2-1\relax}}

% Observaciones y erratas del video. Usar SOLO si hay algo que señalar:
%   \begin{erratas}
%     \errata{203}{3:23}{Se dice ..., pero debe decir ...}
%   \end{erratas}
% El primer argumento son los segundos (enlaza al minuto exacto del video),
% el segundo es la marca visible, el tercero es el texto.
\newenvironment{erratas}{\section*{Observaciones y erratas del video}\begin{itemize}}{\end{itemize}}
\newcommand{\errata}[3]{%
  \edef\errurl{https://www.youtube.com/watch?v=\videoid\string&t=#1}%
  \item[\textbf{\expandafter\href\expandafter{\errurl}{#2}}] #3}
) hace que las
% definiciones de la nota se llamen u.n, u.(n+1), ...; n es el primer número
% libre de esa unidad (ver la tabla de CONVENCIONES.md).
\newcommand{\unidad}[2]{\renewcommand{\thedefn}{#1.\arabic{defn}}\setcounter{defn}{\numexpr#2-1\relax}}

% Observaciones y erratas del video. Usar SOLO si hay algo que señalar:
%   \begin{erratas}
%     \errata{203}{3:23}{Se dice ..., pero debe decir ...}
%   \end{erratas}
% El primer argumento son los segundos (enlaza al minuto exacto del video),
% el segundo es la marca visible, el tercero es el texto.
\newenvironment{erratas}{\section*{Observaciones y erratas del video}\begin{itemize}}{\end{itemize}}
\newcommand{\errata}[3]{%
  \edef\errurl{https://www.youtube.com/watch?v=\videoid\string&t=#1}%
  \item[\textbf{\expandafter\href\expandafter{\errurl}{#2}}] #3}
.
% Compilar con LuaLaTeX:  lualatex nombre-de-la-nota.tex
\usepackage[spanish,provide=*]{babel}
\usepackage[T1]{fontenc}
\usepackage{lmodern}
\usepackage[a4paper,margin=2.6cm]{geometry}
\usepackage{amsmath,amssymb,amsthm}
\usepackage{xcolor}
\usepackage{tikz}
\usetikzlibrary{arrows.meta,babel}
\usepackage[colorlinks=true,urlcolor=blue!50!black,linkcolor=blue!50!black]{hyperref}

\hyphenpenalty=10000 \exhyphenpenalty=10000 \emergencystretch=3em

% Colores de las figuras (los mismos que usa la página)
\definecolor{extremo}{RGB}{47,107,79}
\definecolor{azul}{RGB}{44,74,154}

% Definiciones numeradas por unidad
\theoremstyle{definition}
\newtheorem{defn}{Definición}
% Numeración por unidad: \unidad{u}{n} (después de % Preámbulo compartido de las notas del curso de topología algebraica.
% Cada nota empieza con \documentclass[11pt]{article}, define \videoid{...}
% (el identificador del video de YouTube) y luego hace % Preámbulo compartido de las notas del curso de topología algebraica.
% Cada nota empieza con \documentclass[11pt]{article}, define \videoid{...}
% (el identificador del video de YouTube) y luego hace \input{preambulo}.
% Compilar con LuaLaTeX:  lualatex nombre-de-la-nota.tex
\usepackage[spanish,provide=*]{babel}
\usepackage[T1]{fontenc}
\usepackage{lmodern}
\usepackage[a4paper,margin=2.6cm]{geometry}
\usepackage{amsmath,amssymb,amsthm}
\usepackage{xcolor}
\usepackage{tikz}
\usetikzlibrary{arrows.meta,babel}
\usepackage[colorlinks=true,urlcolor=blue!50!black,linkcolor=blue!50!black]{hyperref}

\hyphenpenalty=10000 \exhyphenpenalty=10000 \emergencystretch=3em

% Colores de las figuras (los mismos que usa la página)
\definecolor{extremo}{RGB}{47,107,79}
\definecolor{azul}{RGB}{44,74,154}

% Definiciones numeradas por unidad
\theoremstyle{definition}
\newtheorem{defn}{Definición}
% Numeración por unidad: \unidad{u}{n} (después de \input{preambulo}) hace que las
% definiciones de la nota se llamen u.n, u.(n+1), ...; n es el primer número
% libre de esa unidad (ver la tabla de CONVENCIONES.md).
\newcommand{\unidad}[2]{\renewcommand{\thedefn}{#1.\arabic{defn}}\setcounter{defn}{\numexpr#2-1\relax}}

% Observaciones y erratas del video. Usar SOLO si hay algo que señalar:
%   \begin{erratas}
%     \errata{203}{3:23}{Se dice ..., pero debe decir ...}
%   \end{erratas}
% El primer argumento son los segundos (enlaza al minuto exacto del video),
% el segundo es la marca visible, el tercero es el texto.
\newenvironment{erratas}{\section*{Observaciones y erratas del video}\begin{itemize}}{\end{itemize}}
\newcommand{\errata}[3]{%
  \edef\errurl{https://www.youtube.com/watch?v=\videoid\string&t=#1}%
  \item[\textbf{\expandafter\href\expandafter{\errurl}{#2}}] #3}
.
% Compilar con LuaLaTeX:  lualatex nombre-de-la-nota.tex
\usepackage[spanish,provide=*]{babel}
\usepackage[T1]{fontenc}
\usepackage{lmodern}
\usepackage[a4paper,margin=2.6cm]{geometry}
\usepackage{amsmath,amssymb,amsthm}
\usepackage{xcolor}
\usepackage{tikz}
\usetikzlibrary{arrows.meta,babel}
\usepackage[colorlinks=true,urlcolor=blue!50!black,linkcolor=blue!50!black]{hyperref}

\hyphenpenalty=10000 \exhyphenpenalty=10000 \emergencystretch=3em

% Colores de las figuras (los mismos que usa la página)
\definecolor{extremo}{RGB}{47,107,79}
\definecolor{azul}{RGB}{44,74,154}

% Definiciones numeradas por unidad
\theoremstyle{definition}
\newtheorem{defn}{Definición}
% Numeración por unidad: \unidad{u}{n} (después de % Preámbulo compartido de las notas del curso de topología algebraica.
% Cada nota empieza con \documentclass[11pt]{article}, define \videoid{...}
% (el identificador del video de YouTube) y luego hace \input{preambulo}.
% Compilar con LuaLaTeX:  lualatex nombre-de-la-nota.tex
\usepackage[spanish,provide=*]{babel}
\usepackage[T1]{fontenc}
\usepackage{lmodern}
\usepackage[a4paper,margin=2.6cm]{geometry}
\usepackage{amsmath,amssymb,amsthm}
\usepackage{xcolor}
\usepackage{tikz}
\usetikzlibrary{arrows.meta,babel}
\usepackage[colorlinks=true,urlcolor=blue!50!black,linkcolor=blue!50!black]{hyperref}

\hyphenpenalty=10000 \exhyphenpenalty=10000 \emergencystretch=3em

% Colores de las figuras (los mismos que usa la página)
\definecolor{extremo}{RGB}{47,107,79}
\definecolor{azul}{RGB}{44,74,154}

% Definiciones numeradas por unidad
\theoremstyle{definition}
\newtheorem{defn}{Definición}
% Numeración por unidad: \unidad{u}{n} (después de \input{preambulo}) hace que las
% definiciones de la nota se llamen u.n, u.(n+1), ...; n es el primer número
% libre de esa unidad (ver la tabla de CONVENCIONES.md).
\newcommand{\unidad}[2]{\renewcommand{\thedefn}{#1.\arabic{defn}}\setcounter{defn}{\numexpr#2-1\relax}}

% Observaciones y erratas del video. Usar SOLO si hay algo que señalar:
%   \begin{erratas}
%     \errata{203}{3:23}{Se dice ..., pero debe decir ...}
%   \end{erratas}
% El primer argumento son los segundos (enlaza al minuto exacto del video),
% el segundo es la marca visible, el tercero es el texto.
\newenvironment{erratas}{\section*{Observaciones y erratas del video}\begin{itemize}}{\end{itemize}}
\newcommand{\errata}[3]{%
  \edef\errurl{https://www.youtube.com/watch?v=\videoid\string&t=#1}%
  \item[\textbf{\expandafter\href\expandafter{\errurl}{#2}}] #3}
) hace que las
% definiciones de la nota se llamen u.n, u.(n+1), ...; n es el primer número
% libre de esa unidad (ver la tabla de CONVENCIONES.md).
\newcommand{\unidad}[2]{\renewcommand{\thedefn}{#1.\arabic{defn}}\setcounter{defn}{\numexpr#2-1\relax}}

% Observaciones y erratas del video. Usar SOLO si hay algo que señalar:
%   \begin{erratas}
%     \errata{203}{3:23}{Se dice ..., pero debe decir ...}
%   \end{erratas}
% El primer argumento son los segundos (enlaza al minuto exacto del video),
% el segundo es la marca visible, el tercero es el texto.
\newenvironment{erratas}{\section*{Observaciones y erratas del video}\begin{itemize}}{\end{itemize}}
\newcommand{\errata}[3]{%
  \edef\errurl{https://www.youtube.com/watch?v=\videoid\string&t=#1}%
  \item[\textbf{\expandafter\href\expandafter{\errurl}{#2}}] #3}
) hace que las
% definiciones de la nota se llamen u.n, u.(n+1), ...; n es el primer número
% libre de esa unidad (ver la tabla de CONVENCIONES.md).
\newcommand{\unidad}[2]{\renewcommand{\thedefn}{#1.\arabic{defn}}\setcounter{defn}{\numexpr#2-1\relax}}

% Observaciones y erratas del video. Usar SOLO si hay algo que señalar:
%   \begin{erratas}
%     \errata{203}{3:23}{Se dice ..., pero debe decir ...}
%   \end{erratas}
% El primer argumento son los segundos (enlaza al minuto exacto del video),
% el segundo es la marca visible, el tercero es el texto.
\newenvironment{erratas}{\section*{Observaciones y erratas del video}\begin{itemize}}{\end{itemize}}
\newcommand{\errata}[3]{%
  \edef\errurl{https://www.youtube.com/watch?v=\videoid\string&t=#1}%
  \item[\textbf{\expandafter\href\expandafter{\errurl}{#2}}] #3}

\unidad{1}{5} % unidad 1; primera definición: 1.5

\title{La homotopía de caminos es una relación de equivalencia}
\author{Luis Jorge Sánchez Saldaña\\[2pt]\small Facultad de Ciencias, UNAM}
\date{}

\begin{document}
\maketitle

\noindent\small Notas del video
\href{https://www.youtube.com/watch?v=qFL_OneuFwI}{\guillemotleft{}La homotopía de caminos rel 0,1 es una relación de equivalencia\guillemotright{}}
(curso de Introducción a la Topología Algebraica). Demuestran que ser homotópicos relativo a los extremos es una relación de equivalencia en el conjunto de caminos.
\normalsize

\section*{El espacio de caminos y la relación}

Sea $X$ un espacio topológico cualquiera, fijo durante todo el video. Se recuerda que un camino es una función continua del intervalo unitario $I$ en $X$.

\begin{defn}[Espacio de caminos y la relación $\sim$]
Se denota por $\Omega(X)$ al conjunto de todos los caminos en $X$,
\[
  \Omega(X)=\{\,f\colon I\to X\ \mid\ f \text{ es continua}\,\}.
\]
Para $f_1,f_2\in\Omega(X)$ se dice que $f_1$ está relacionado con $f_2$, y se escribe $f_1\sim f_2$, si
\[
  f_1\simeq f_2\ \ \mathrm{rel}\ \{0,1\},
\]
es decir, si son homotópicos relativos a los extremos.
\end{defn}

El teorema principal del video es que esta relación es en realidad una relación de equivalencia.

\begin{prop}
La relación $\sim$ en $\Omega(X)$ es una relación de equivalencia.
\end{prop}

Para demostrarlo hay que probar que es reflexiva, simétrica y transitiva. Se procede por partes.

\begin{proof}[Reflexividad]
Se quiere ver que cualquier camino $f$ en $X$ cumple $f\simeq f\ \mathrm{rel}\ \{0,1\}$. Intuitivamente, un camino se puede deformar a sí mismo dejando fijos los extremos sin deformarlo: basta tomar la deformación que es el camino quieto todo el tiempo. Formalmente, la homotopía $F\colon I\times I\to X$ dada por
\[
  F(s,t)=f(s)
\]
manda cada horizontal del cuadrado (cada $t$ fijo) al mismo camino $f$. Además, los extremos siempre quedan fijos, pues $F(0,t)=f(0)$ y $F(1,t)=f(1)$ para todo $t$. Así, $F$ es una homotopía de $f$ en $f$ relativa a $\{0,1\}$.
\end{proof}

\begin{proof}[Simetría]
Se quiere ver que si $f\simeq g\ \mathrm{rel}\ \{0,1\}$, entonces $g\simeq f\ \mathrm{rel}\ \{0,1\}$.

Antes de la demostración formal, la idea: si se puede deformar $f$ en $g$ dejando fijos los extremos, también se puede empezar con $g$ y deformar en $f$; basta tomar esa deformación pero al revés. En el cuadrado $I\times I$, esto es reflejarlo: la homotopía que se tiene empieza con $f$ arriba y va deformando hasta $g$ abajo, y al reflejar el cuadrado se empieza con $g$ y se termina con $f$, y los extremos siguen fijos aunque se vaya de abajo hacia arriba.

Formalmente, sea $F\colon I\times I\to X$ una homotopía de $f$ a $g$ relativa a los extremos. Se define $G\colon I\times I\to X$ como
\[
  G(s,t)=F(s,1-t),
\]
es decir, se cambia $t$ por $1-t$, lo que \guillemotleft{}voltea\guillemotright{} la homotopía. El video deja como un ejercicio muy bonito verificar que $G$ es, en efecto, una homotopía de $g$ a $f$ relativa a los extremos.
\end{proof}

\begin{proof}[Transitividad]
Se quiere ver que si $f\simeq g\ \mathrm{rel}\ \{0,1\}$ y $g\simeq h\ \mathrm{rel}\ \{0,1\}$, entonces $f\simeq h\ \mathrm{rel}\ \{0,1\}$.

La idea: se puede deformar $f$ en $g$ dejando fijos los extremos y luego deformar $g$ en $h$ dejando fijos los extremos; basta pegar las dos deformaciones, primero deformando $f$ hasta llegar a $g$ y luego siguiendo de largo hasta $h$. Esquemáticamente, se pegan los dos cuadrados de las homotopías, uno encima del otro (figura). El único detalle es que al pegarlos se obtiene un rectángulo de tamaño $2\times1$ y no el cuadrado unitario, por lo que hay que reescalar para que la nueva homotopía salga de $I\times I$.

\begin{figure}[ht]
\centering
\begin{tikzpicture}[>=Stealth,line cap=round]
  \def\w{3.2}\def\hh{3.2}
  \draw[very thick] (0,{\hh}) -- (\w,{\hh});
  \draw[thick,gray] (0,{\hh/2}) -- (\w,{\hh/2});
  \draw[very thick,azul] (0,0) -- (\w,0);
  \draw[very thick,extremo] (0,0) -- (0,{\hh});
  \draw[very thick,extremo] (\w,0) -- (\w,{\hh});
  \node at ({\w/2},{3*\hh/4}) {$F$};
  \node at ({\w/2},{\hh/4}) {$G$};
  \node[left] at (0,{\hh}) {$t=0\mapsto f$};
  \node[left] at (0,{\hh/2}) {$t=\tfrac12\mapsto g$};
  \node[left,azul] at (0,0) {$t=1\mapsto h$};
  \node[gray,below] at ({\w/2},-0.45) {\small $I\times I$};
\end{tikzpicture}
\par\smallskip
\begin{minipage}{0.9\textwidth}\small\centering
La parte de arriba es la homotopía $F$ de $f$ a $g$ y la de abajo es la homotopía $G$ de $g$ a $h$; los lados verticales (en verde) van a los extremos.
\end{minipage}
\end{figure}

Formalmente, sean $F\colon I\times I\to X$ y $G\colon I\times I\to X$ homotopías de $f$ a $g$ y de $g$ a $h$, respectivamente, ambas relativas a los extremos. Se define $H\colon I\times I\to X$ por casos como
\[
  H(s,t)=
  \begin{cases}
    F(s,2t) & \text{si } t\in[0,\tfrac12],\\[2pt]
    G(s,2t-1) & \text{si } t\in[\tfrac12,1].
  \end{cases}
\]
El video deja como un muy buen ejercicio verificar que $H$ es una homotopía de $f$ a $h$ relativa a los extremos: hay que revisar con cuidado las definiciones y evaluar.
\end{proof}

\section*{Conclusión}

Con esto queda demostrado que la relación de homotopía relativa a los extremos es una relación de equivalencia en el conjunto de todos los caminos de $X$. Este es un paso más hacia la definición del grupo fundamental.

\begin{erratas}
  \errata{110}{1:50}{En la reflexividad se escribe $F(s,t)=f(t)$, pero debe decir $F(s,t)=f(s)$ (con $f(t)$ los extremos no quedan fijos).}
  \errata{255}{4:15}{En la simetría se define $G(s,t)=F(1-s,t)$, pero debe decir $G(s,t)=F(s,1-t)$ (así $G$ va de $g$ a $f$; con $F(1-s,t)$ solo se recorre cada camino al revés).}
\end{erratas}

\end{document}
